\begin{ejercicio}
    % ID: 19
    Mostrar que $\Phi(x)$ es un lazo y que $\Phi$ es biyectiva.
\end{ejercicio}

\begin{demostracion}
    Recordemos que $\Phi_1:\map(X\times I,Y)\to\map(X,\map(I,Y))$ induce la función 
    \begin{equation*}
    \Phi:\map((SX,*),(Y,y_0))\to \map((X,x_0),(\Omega(Y,y_0),c_{y_0}))
\end{equation*}
tal que $\Phi(f)=\Phi_1(f\circ p)$.
    
    Queremos mostrar que $\Phi_1(f\circ p)(x)$ es un lazo basada en $y_0$ para toda $x\in X$. Por definición
    \begin{align*}
        \Phi(f\circ p)(x)(t)&=(f\circ p)(x,t)\\
        &=f[x,t]
    \end{align*}
    Como $(x,0),(x,1)\in (X\times \partial I)\cup (\{x_0\}\times I)$, entonces $[x,0]=[x,1]=*$, como $f$ es basada, se sigue que
    \begin{equation*}
        f[x,0]=f[x,1]=f(*)=y_0.
    \end{equation*}
    Finalmente, como $\Phi_1(f\circ p)(x)$ es continua, se sigue que $\Phi_1(f \circ p)(x)$ es un lazo. Además 
    \begin{equation*}
        \Phi(f\circ p)(x_0)(t)=f[x_0,t]=f(*)=y_0
    \end{equation*}
    por lo tanto, $\Phi_1(f\circ p)(x_0)=c_{y_0}$, i.e., $\Phi$ está bien definida.
    
    Mostremos ahora que $\Phi$ es una biyección. Tomemos 
    $$g\in \map((X,x_0),(\Omega(Y,y_0),c_{y_0}))\subseteq \map(X,\map(I, Y))$$
    Como $I$ localmente compacto Hausdorff, existe la inversa de $\Phi_1$, entonces $G=\Phi^{-1}_1(g):X\times I\to Y$ es continua tal que $G(x,t)=g(x)(t)$ (por definición de $\Phi_1^{-1}$). Sea $(x,t)\in X\times \partial I$, como $g(x)\in \Omega(Y,y_0)$, entonces
    \begin{equation*}
        G(x,t)=g(x)(t)=y_0
    \end{equation*}
    Por otro lado,como $g(x_0)=c_{y_0}$, entonces
        \begin{equation*}
        G(x_0,t)=g(x_0)(t)=y_0, \text{ para todo } t\in I.
    \end{equation*}
    entonces $G$ pasa al cociente, i.e., existe una única función continua $\overline{G}:SX\to Y$ tal que $\overline{G}\circ p=G$. Mostremos que $\overline{G}$ es basado en $*$. Notemos que $p(x_0,t)=*$, entonces
    \begin{equation*}
        \overline{G}(*)=\overline{G}(p(x_0,t))=G(x_0,t)=y_0.
    \end{equation*}
    Por lo anterior, podemos definir  
    \begin{align*}
        \Psi:\map((X,x_0),(\Omega(Y,y_0),c_{y_0}))&\to \map((SX,*),(Y,y_0))\\
        g&\mapsto \overline{G}=\overline{\Phi_1^{-1}(g)}
    \end{align*}
    Mostremos que $\Psi$ es la inversa de $\Phi$. Por definición $\Psi(\Phi(f))$ es la única función continua $\overline{F}:SX\to Y$ tal que $\overline{F}\circ p=F$, donde $F=\Phi_1^{-1}(\Phi_1(f\circ p))=f\circ p$, por la unicidad de $\overline{F}$, se sigue que 
    \begin{equation*}
        \Psi(\Phi(f))=\overline{F}=f
    \end{equation*}
    Por otro lado,
    \begin{align*}
        \Phi(\Psi(g))&=\Phi_1(\Psi(g)\circ p)\\
        &=\Phi_1(\overline{G}\circ p)\\
        &=\Phi_1(G)\\
        &=\Phi_1(\Phi_1^{-1}(g))\\
        &=g.
    \end{align*}
    Lo anterior muestra que $\Phi$ es una biyección con inversa $\Psi$.
\end{demostracion}
